\documentclass[aspectratio=169]{beamer}

\usepackage[utf8]{inputenc}
\usepackage[T1]{fontenc}

\usepackage[sfdefault,light]{FiraSans}
\renewcommand{\familydefault}{\sfdefault}
\usefonttheme{professionalfonts}

\usepackage{xcolor}
\usepackage{graphicx}
\usepackage{ifthen}
\usepackage{hyperref}
\hypersetup{hidelinks}
\usepackage[strings]{underscore}

\usetheme{default}

% =====================================================
% Brand (MDA / Núcleo de Inteligência Artificial)
% Files live in: figs/
% =====================================================
\newcommand{\logoSmall}{mda.jpg}
\newcommand{\logoBig}{semeia.png}
\newcommand{\brandshort}{META-3 | SEMEIA}

% =====================================================
% Placeholders (edit these)
% =====================================================
\newcommand{\DeckTitle}{\textbf{SEMEIA Systems Engineering: Briefing \& Roadmap}}
\newcommand{\DeckSubtitle}{} % set to empty {} if you want none
\newcommand{\PresenterName}{META-3 Coordination Prof. Bruno}
\newcommand{\PresenterEmail}{Automation Tech lead: Antonio}
\newcommand{\OrgLine}{AI Tech Lead: Mohamed}
\newcommand{\DateLine}{}

% Colors
\definecolor{mdagrey}{RGB}{90,90,90}
\definecolor{mdalightgrey}{RGB}{200,200,200}
\definecolor{accentorange}{RGB}{230,90,40}

% Base text and titles
\setbeamerfont{normal text}{size=8pt}
\setbeamerfont{frametitle}{size=\normalsize,series=\bfseries}
\setbeamerfont{title}{size=\Large,series=\bfseries}
\setbeamercolor{title}{fg=accentorange}
\setbeamercolor{frametitle}{fg=black}
\setbeamercolor{normal text}{fg=black}

% Compact bullets
\setbeamerfont{itemize/enumerate body}{size=\scriptsize}
\setbeamertemplate{itemize item}{}
\setlength{\leftmargini}{1.2em}
\setlength{\itemsep}{0.28em}

% Margins
\setbeamersize{text margin left=1.35cm, text margin right=1.35cm}
\setlength{\columnsep}{0.9cm}

% =====================================================
% Graphics macro (expands macros like \logoBig)
% =====================================================
\makeatletter
\newcommand{\deckgraphic}[2][]{%
  \begingroup
  \edef\dg@file{#2}%
  \IfFileExists{figs/\dg@file}{%
    \includegraphics[
      width=\linewidth,
      height=0.82\textheight,
      keepaspectratio,
      #1
    ]{figs/\dg@file}%
  }{%
    {\color{red}\fbox{\scriptsize Missing: figs/\dg@file}}%
  }%
  \endgroup
}
\makeatother

\newcommand{\focusline}[1]{%
  \vspace{0.25em}%
  {\color{accentorange}\footnotesize\bfseries #1}%
}

% Footer reference (per-slide citation footer)
\newcommand{\insertframeref}{}
\newcommand{\frameref}[1]{\gdef\insertframeref{#1}}

\setbeamertemplate{headline}{}
\setbeamertemplate{footline}{%
  \leavevmode%
  \hbox{%
    \begin{beamercolorbox}[wd=\paperwidth,ht=1.6ex,dp=0.5ex,leftskip=0.3cm]{footline}%
      {\color{mdagrey}\fontsize{6pt}{4.8pt}\selectfont\insertframeref}%
    \end{beamercolorbox}%
  }%
}
\setbeamertemplate{navigation symbols}{}

% Header row with small logo at right
\setbeamertemplate{frametitle}{%
  \vspace*{0.05cm}%
  \begin{beamercolorbox}[wd=\paperwidth,ht=2.4ex,dp=0.6ex,leftskip=0.35cm,rightskip=0.35cm]{frametitle}%
    {\fontsize{6pt}{7pt}\selectfont
      \color{mdagrey}%
      \makebox[0.30\paperwidth][l]{\insertframetitle}%
      \makebox[0.40\paperwidth][c]{\brandshort}%
      \makebox[0.30\paperwidth][c]{\hspace*{-0.1cm}\deckgraphic[height=0.5cm]{\logoSmall}}%
    }%
  \end{beamercolorbox}%
  \vspace{0.15cm}%
}

% Two-column slide macro
\newcommand{\twocolslide}[5]{%
  \begin{frame}
    \frametitle{#2}%
    \frameref{#1}%
    \begin{columns}[c,onlytextwidth]
      \begin{column}{0.38\textwidth}
        \textbf{#3}
        \vspace{0.3em}
        #4
      \end{column}
      \begin{column}{0.58\textwidth}
        #5
      \end{column}
    \end{columns}
  \end{frame}
}

% =====================================================
% Outline separator macro (5 sections total)
% =====================================================
\newcommand{\outlineseparator}[1]{%
  \begin{frame}
    \frametitle{Outline}%
    \frameref{}%
    \centering
    \vfill
    {\Large
      \ifnum#1=1 \textbf{SEMEIA Overview (incl.\ UX)}\else {\color{mdalightgrey}SEMEIA Overview (incl.\ UX)}\fi \\[0.55em]
      \ifnum#1=2 \textbf{SEMEIA UI Gateway}\else {\color{mdalightgrey}SEMEIA UI Gateway}\fi \\[0.55em]
      \ifnum#1=3 \textbf{SEMEIA Backend AI System}\else {\color{mdalightgrey}SEMEIA Backend AI System}\fi \\[0.55em]
      \ifnum#1=4 \textbf{SEMEIA Integration Gateway}\else {\color{mdalightgrey}SEMEIA Integration Gateway}\fi \\[0.55em]
      \ifnum#1=5 \textbf{SEMEIA Evolution}\else {\color{mdalightgrey}SEMEIA Evolution}\fi
    }
    \vfill
  \end{frame}
}

% Title info
\title[\DeckTitle]{\DeckTitle}
\institute{%
  \texttt{\PresenterName}\\[0.15em]
  \texttt{\PresenterEmail}\\[0.15em]
  \texttt{\OrgLine}\\[0.15em]
  \texttt{\DateLine}\\[0.15em]
}
\date{}

\begin{document}

% -----------------------------------------------------
% Slide 1 – Title
% -----------------------------------------------------
\begin{frame}[plain]
  \frameref{}%
  \begin{columns}[T,onlytextwidth]
    \begin{column}{0.62\textwidth}
      \vspace{0.85cm}
      {\LARGE\bfseries \DeckTitle\par}

      \ifthenelse{\equal{\DeckSubtitle}{}}{}{%
        \vspace{0.30cm}
        {\large \DeckSubtitle\par}
      }

      \vspace{0.85cm}
      {\small \insertinstitute\par}
    \end{column}

    \begin{column}{0.38\textwidth}
      \vspace{0.20cm}
      \hfill
      \deckgraphic[height=0.86\textheight]{\logoBig}
    \end{column}
  \end{columns}
\end{frame}

% -----------------------------------------------------
% Section 1
% -----------------------------------------------------
\outlineseparator{1}

% =====================================================
% 1) SEMEIA Overview (incl. UX)
% =====================================================

\begin{frame}
  \frametitle{SEMEIA Overview}
  \frameref{}%
  \centering
  \IfFileExists{figs/overview.png}{
    \includegraphics[
      width=0.95\textwidth,
      height=0.85\textheight,
      keepaspectratio
    ]{figs/overview.png}
  }{
    {\color{mdagrey}\scriptsize Add figs/overview.png (optional) to display the overview poster.}
  }
\end{frame}

\twocolslide
  {}
  {SEMEIA Overview}
  {Mission}
  {
    \begin{itemize}
      \item Reduce operational burden via automation
      \item Improve service delivery quality
      \item Standardize workflows across systems
    \end{itemize}
  }
  {
    \focusline{SEMEIA is a systems program: UX + AI + integration + governance.}
  }

% --- Conversational AI framing
\twocolslide
  {%
    1.~Atlassian (Conversational AI) \href{https://www.atlassian.com/blog/artificial-intelligence/conversation-ai}{atlassian.com/blog/artificial-intelligence/conversation-ai}
  }
  {SEMEIA Overview}
  {SEMEIA as Conversational AI}
  {
    \begin{itemize}
      \item SEMEIA is a \textbf{conversational AI system}\textsuperscript{1}
      \item Primary interface: \textbf{chat} (natural language)
      \item Executes workflows through controlled gateways
    \end{itemize}
    \vfill
    \focusline{Conversation = interface + workflow engine.}
  }
  {
    \begin{itemize}
      \item Supports:
        \begin{itemize}
          \item multi-turn context
          \item clarification loops
          \item structured outputs (forms, drafts)
          \item tool execution with auditability
        \end{itemize}
    \end{itemize}
  }

% --- REST vs GraphQL
\twocolslide
  {%
    1.~Fielding (2000) REST dissertation \href{https://www.ics.uci.edu/~fielding/pubs/dissertation/rest_arch_style.htm}{ics.uci.edu/\textasciitilde fielding/...};\;
    2.~GraphQL Spec \href{https://spec.graphql.org/}{spec.graphql.org}
  }
  {SEMEIA Overview}
  {REST vs.\ GraphQL (Contract Strategy)}
  {
    \begin{itemize}
      \item \textbf{REST}\textsuperscript{1}: stable resources + caching
      \item \textbf{GraphQL}\textsuperscript{2}: UI/agent-friendly query surface
      \item SEMEIA uses \textbf{contracts} as governance
    \end{itemize}
    \vfill
    \focusline{GraphQL is the interface; persisted ops are the guardrail.}
  }
  {
    \begin{itemize}
      \item Design target:
        \begin{itemize}
          \item REST-like stable execution endpoints
          \item GraphQL for orchestration / data shaping
          \item no ad-hoc upstream access by agents
        \end{itemize}
    \end{itemize}
  }

% --- LLM history
\twocolslide
  {%
    1.~Vaswani et al.\ (2017) Transformers \href{https://arxiv.org/abs/1706.03762}{arxiv.org/abs/1706.03762}
  }
  {SEMEIA Overview}
  {LLMs History (High-level)}
  {
    \begin{itemize}
      \item Transformers enabled scalable attention\textsuperscript{1}
      \item Large-scale pretraining $\rightarrow$ general language competence
      \item Instruction tuning $\rightarrow$ usable assistants
    \end{itemize}
    \vfill
    \focusline{LLMs changed UX: natural language became a control plane.}
  }
  {
    \begin{itemize}
      \item SEMEIA impact:
        \begin{itemize}
          \item drafting automation
          \item information extraction
          \item orchestration via tools
        \end{itemize}
    \end{itemize}
  }

% --- Agents / learning framing
\twocolslide
  {}
  {SEMEIA Overview}
  {Autonomous Agents vs.\ Learning Agents}
  {
    \begin{itemize}
      \item \textbf{Autonomous agent}: plans + executes via tools
      \item \textbf{Learning agent}: improves via feedback + data
      \item Public sector requires \textbf{HITL approval}
    \end{itemize}
    \vfill
    \focusline{Autonomy without governance is operational risk.}
  }
  {
    \begin{itemize}
      \item SEMEIA approach:
        \begin{itemize}
          \item controlled action space (persisted ops)
          \item audit logs and explainability
          \item human approvals for sensitive actions
        \end{itemize}
    \end{itemize}
  }

% --- MCP vs Skills
\twocolslide
  {%
    1.~Block/Goose (2025-12-22) \href{https://block.github.io/goose/blog/2025/12/22/agent-skills-vs-mcp/}{block.github.io/goose/...}
  }
  {SEMEIA Overview}
  {MCP Server vs.\ Skills for Agents}
  {
    \begin{itemize}
      \item \textbf{MCP}\textsuperscript{1}: capabilities layer (tools)
      \item \textbf{Skills}\textsuperscript{1}: behavior layer (how to use tools)
    \end{itemize}
    \vfill
    \focusline{MCP gives abilities; Skills teach safe usage.}
  }
  {
    \begin{itemize}
      \item Key idea:
        \begin{itemize}
          \item skills without MCP = instructions only
          \item MCP without skills = power without guidance
        \end{itemize}
    \end{itemize}
  }

% =====================================================
% UX content INSIDE Overview (improved)
% =====================================================

\begin{frame}
  \frametitle{SEMEIA UX (within Overview)}
  \frameref{}%
  \centering
  \begin{minipage}{0.48\textwidth}
    \centering
    \IfFileExists{figs/ui.png}{
      \includegraphics[width=\linewidth,keepaspectratio]{figs/ui.png}
    }{
      {\color{mdagrey}\scriptsize Add figs/ui.png}
    }
  \end{minipage}
  \hfill
  \begin{minipage}{0.48\textwidth}
    \centering
    \IfFileExists{figs/sei.jpg}{
      \includegraphics[width=\linewidth,keepaspectratio]{figs/sei.jpg}
    }{
      {\color{mdagrey}\scriptsize Add figs/sei.jpg (optional)}
    }
  \end{minipage}
\end{frame}

\twocolslide
  {}
  {SEMEIA Overview}
  {UX Goals (Improved)}
  {
    \begin{itemize}
      \item Reduce cognitive load (user does not learn the bureaucracy)
      \item Reduce context switching (SEI/TransfereGov/etc.)
      \item Convert policies into guided steps
      \item Default to ``ask-back'' for missing inputs
    \end{itemize}
  }
  {
    \begin{itemize}
      \item UX outcomes:
        \begin{itemize}
          \item fewer incorrect requests
          \item faster completion
          \item higher trust in automation
        \end{itemize}
    \end{itemize}
    \vspace{0.3em}
    \focusline{Simplicity is not aesthetics; it is operational efficiency.}
  }

\twocolslide
  {}
  {SEMEIA Overview}
  {UX Architecture Requirements}
  {
    \begin{itemize}
      \item Multi-step workflows with state
      \item Explainability in UI (why this step?)
      \item ``Preview before execute'' for actions
      \item Evidence view (sources, IDs, logs)
    \end{itemize}
  }
  {
    \begin{itemize}
      \item Failure-handling UX:
        \begin{itemize}
          \item retries with reasons
          \item clear responsibility handoffs
          \item partial results are allowed
        \end{itemize}
    \end{itemize}
  }

% --- Overview roadmap
\twocolslide
  {}
  {SEMEIA Overview}
  {Roadmap (Overview)}
  {
    \begin{itemize}
      \item \textbf{Now (0--2 weeks)}
        \begin{itemize}
          \item scope + glossary freeze
          \item UX journey mapping + ownership
          \item integration dependency registry
        \end{itemize}
      \item \textbf{Next (2--6 weeks)}
        \begin{itemize}
          \item architecture blueprint
          \item workflow MVP definitions
        \end{itemize}
    \end{itemize}
  }
  {
    \begin{itemize}
      \item \textbf{Later (6+ weeks)}
        \begin{itemize}
          \item governance playbook + approvals
          \item KPI dashboards
        \end{itemize}
      \item \textbf{Risks}: unclear contracts; cross-system UX ambiguity
    \end{itemize}
  }

% -----------------------------------------------------
% Section 2
% -----------------------------------------------------
\outlineseparator{2}

% =====================================================
% 2) SEMEIA UI Gateway
% =====================================================

\begin{frame}
  \frametitle{SEMEIA UI Gateway (UIGW)}
  \frameref{}%
  \centering
  \vfill
  {\Large\bfseries UI Gateway / BFF Layer}
  \vspace{1em}

  {\small\color{mdagrey}
  UX-friendly API layer for SEMEIA UI + caching + permissions + orchestration}
  \vfill
\end{frame}

\twocolslide
  {}
  {SEMEIA UI Gateway}
  {What is UIGW?}
  {
    \begin{itemize}
      \item UIGW = Backend-for-Frontend (BFF)
      \item Converts workflows into UI-friendly contracts
      \item Aggregates multiple backend calls
    \end{itemize}
  }
  {
    \focusline{UIGW optimizes for UX without exposing provider complexity.}
  }

% -----------------------------------------------------
% Requirements slide (EXACT STYLE requested)
% -----------------------------------------------------
\twocolslide
  {%
    1.~Fielding (2000) REST dissertation \href{https://www.ics.uci.edu/~fielding/pubs/dissertation/rest_arch_style.htm}{ics.uci.edu/\textasciitilde fielding/...};\;
    2.~Nygard (2007) \textit{Release It!};\;
    3.~Google SRE Book (2016) \href{https://sre.google/sre-book/table-of-contents/}{sre.google/sre-book}
  }
  {SEMEIA UI Gateway}
  {Operational \& Architectural Requirements}
  {
    \begin{itemize}
      \item \textbf{Design principles:} stable screen contracts + separation (UI vs integration)\textsuperscript{1}
      \item \textbf{Why:} protect UX from backend volatility and failures\textsuperscript{2}
    \end{itemize}
    \vfill
    \focusline{UIGW is the UX reliability boundary.}
  }
  {
    \centering
    \begin{itemize}
      \item \textbf{R1} – Screen-oriented APIs (BFF): \textit{what the UI needs}
      \item \textbf{R2} – RBAC enforcement aligned with UI actions
      \item \textbf{R3} – Response shaping: stable view models
      \item \textbf{R4} – Caching policy (TTL per endpoint; per-user keys)
      \item \textbf{R5} – Rate limiting + payload caps (UI safety)
      \item \textbf{R6} – Degrade gracefully (partial responses allowed)\textsuperscript{3}
      \item \textbf{R7} – Observability: request\_id, structured logs, per-route metrics\textsuperscript{3}
    \end{itemize}
  }

\twocolslide
  {}
  {SEMEIA UI Gateway}
  {Roadmap (UIGW)}
  {
    \begin{itemize}
      \item \textbf{Now (0--2 weeks)}: define screen contracts, error model
      \item \textbf{Next (2--6 weeks)}: implement dashboard/task APIs + caching
      \item \textbf{Later (6+ weeks)}: performance tuning + incident tooling
    \end{itemize}
  }
  {
    \begin{itemize}
      \item Owners: UI backend engineers + Tech Lead
      \item Risks: RBAC mismatches; UI scope drift
    \end{itemize}
  }

% -----------------------------------------------------
% Section 3
% -----------------------------------------------------
\outlineseparator{3}

% =====================================================
% 3) SEMEIA Backend AI System
% =====================================================

\begin{frame}
  \frametitle{SEMEIA Backend AI System}
  \frameref{}%
  \centering
  \vfill
  {\Large\bfseries AI Services + Orchestration + Governance}
  \vspace{1em}

  {\small\color{mdagrey}
  Models, prompts, tools, evaluation, monitoring}
  \vfill
\end{frame}

\twocolslide
  {}
  {SEMEIA Backend AI}
  {Core Components}
  {
    \begin{itemize}
      \item Orchestrator (agent runtime)
      \item Skill library + tools
      \item Retrieval (KB/search)
      \item Evaluation harness
    \end{itemize}
  }
  {
    \focusline{Backend AI is a system, not a model.}
  }

\twocolslide
  {%
    1.~Manning: \textit{Build a Large Language Model (From Scratch)};\;
    2.~Manning: \textit{Build a Reasoning Large Language Model (From Scratch)}
  }
  {SEMEIA Backend AI}
  {LLMs vs.\ Reasoning LLMs}
  {
    \begin{itemize}
      \item \textbf{LLMs}\textsuperscript{1}: drafting, summarizing, extraction
      \item \textbf{Reasoning LLMs}\textsuperscript{2}: planning, decomposition, multi-step workflows
      \item SEMEIA requires both: conversation + execution correctness
    \end{itemize}
    \vfill
    \focusline{Reasoning must be engineered: tools, checks, and guardrails.}
  }
  {
    \begin{itemize}
      \item Impact on SEMEIA:
        \begin{itemize}
          \item workflow completion reliability
          \item fewer hallucinated parameters
          \item stronger tool orchestration
        \end{itemize}
    \end{itemize}
  }

% --- Mitchell quote slide
\twocolslide
  {%
    1.~Tom M. Mitchell (1997) \textit{Machine Learning}
  }
  {SEMEIA Backend AI}
  {What is Machine Learning? (Mitchell)}
  {
    \begin{itemize}
      \item Learning loops require:
        \begin{itemize}
          \item feedback collection
          \item test sets and regression checks
          \item measurable improvement criteria
        \end{itemize}
    \end{itemize}
    \vfill
    \focusline{Learning is a lifecycle, not a feature.}
  }
  {
    \vfill
    {\footnotesize
    ``The field of machine learning is concerned with the question of how to construct computer programs that automatically improve with experience.''\textsuperscript{1}}
    \vfill
  }

% -----------------------------------------------------
% Requirements slide (EXACT STYLE requested)
% -----------------------------------------------------
\twocolslide
  {%
    1.~Mitchell (1997) \textit{Machine Learning};\;
    2.~Vaswani et al.\ (2017) Transformers \href{https://arxiv.org/abs/1706.03762}{arxiv.org/abs/1706.03762};\;
    3.~Google SRE Book (2016) \href{https://sre.google/sre-book/table-of-contents/}{sre.google/sre-book}
  }
  {SEMEIA Backend AI System}
  {Operational \& Architectural Requirements}
  {
    \begin{itemize}
      \item \textbf{Design principles:} measurable learning + reliability via engineered pipelines\textsuperscript{1}
      \item \textbf{Why:} LLM behavior changes; system must remain stable\textsuperscript{2}
    \end{itemize}
    \vfill
    \focusline{The AI system must be reproducible, observable, and governable.}
  }
  {
    \centering
    \begin{itemize}
      \item \textbf{R1} – Controlled action space (persisted ops only)
      \item \textbf{R2} – Human-in-the-loop approvals for sensitive actions
      \item \textbf{R3} – Evaluation harness: regression tests for prompts/tools
      \item \textbf{R4} – Data governance: sources, freshness, provenance
      \item \textbf{R5} – Prompt/skill versioning + change management
      \item \textbf{R6} – Observability: traces, tool call logs, cost/latency metrics\textsuperscript{3}
      \item \textbf{R7} – Safe failure modes (ask-back, stop, escalate)
    \end{itemize}
  }

\twocolslide
  {}
  {SEMEIA Backend AI}
  {Roadmap (AI System)}
  {
    \begin{itemize}
      \item \textbf{Now}: define skill library + tool registry + eval baseline
      \item \textbf{Next}: implement top workflows + approvals
      \item \textbf{Later}: drift monitoring + full regression suite
    \end{itemize}
  }
  {
    \begin{itemize}
      \item Owners: AI Eng + MLOps + Tech Lead
      \item Risks: unclear approval boundaries; tool misuse
    \end{itemize}
  }

% -----------------------------------------------------
% Section 4
% -----------------------------------------------------
\outlineseparator{4}

% =====================================================
% 4) SEMEIA Integration Gateway
% =====================================================

\begin{frame}
  \frametitle{SEMEIA Integration Gateway Requirements}
  \frameref{}%
  \centering
  \vfill
  {\Large\bfseries Operational \& Architectural Requirements}
  \vspace{1em}

  {\small\color{mdagrey}
  Gateway + Adapters + Persisted Ops + Observability + Security}
  \vfill
\end{frame}

\twocolslide
  {}
  {SEMEIA Integration Gateway (IGW)}
  {What is IGW?}
  {
    \begin{itemize}
      \item Policy boundary between agents and external systems
      \item Executes \textbf{persisted operations} only
      \item Standard error model + schema validation
      \item Auditable calls (IDs, logs, metrics)
    \end{itemize}
  }
  {
    \focusline{IGW prevents ad-hoc API access by agents.}
  }

% -----------------------------------------------------
% Requirements slide (EXACT STYLE requested)
% -----------------------------------------------------
\twocolslide
  {%
    1.~Nygard (2007) \textit{Release It!};\;
    2.~Google SRE Book (2016) \href{https://sre.google/sre-book/table-of-contents/}{sre.google/sre-book};\;
    3.~GraphQL Spec \href{https://spec.graphql.org/}{spec.graphql.org}
  }
  {SEMEIA Integration Gateway (IGW)}
  {Operational \& Architectural Requirements}
  {
    \begin{itemize}
      \item \textbf{Design principles:} guardrails-first integration boundary\textsuperscript{3}
      \item \textbf{Why:} integration failures are inevitable; system must isolate faults\textsuperscript{1}
    \end{itemize}
    \vfill
    \focusline{Persisted ops are the governance primitive of SEMEIA integrations.}
  }
  {
    \centering
    \begin{itemize}
      \item \textbf{R1} – Persisted ops only (no arbitrary upstream calls)
      \item \textbf{R2} – Config-driven routing + strict versioning
      \item \textbf{R3} – Deterministic schema validation (JSONSchema)
      \item \textbf{R4} – Standard error model + retry policy classes
      \item \textbf{R5} – Fault isolation (adapter failures never crash gateway)\textsuperscript{1}
      \item \textbf{R6} – Observability: request\_id everywhere + metrics\textsuperscript{2}
      \item \textbf{R7} – Security: secrets never in specs; least privilege
    \end{itemize}
  }

\twocolslide
  {}
  {SEMEIA Integration Gateway (IGW)}
  {Roadmap (IGW)}
  {
    \begin{itemize}
      \item \textbf{Now}: spec freeze + skeleton + logs
      \item \textbf{Next}: onboard providers + schema enforcement
      \item \textbf{Later}: circuit breaker + governance CI
    \end{itemize}
  }
  {
    \begin{itemize}
      \item Owners: Platform Eng + Integration Eng
      \item Risks: upstream contract drift; ops explosion
    \end{itemize}
  }

% -----------------------------------------------------
% Section 5
% -----------------------------------------------------
\outlineseparator{5}

% =====================================================
% 5) SEMEIA Evolution
% =====================================================

\twocolslide
  {}
  {SEMEIA Evolution}
  {Incremental Development}
  {
    \begin{itemize}
      \item Iterative rollout (MVP $\rightarrow$ expansions)
      \item Governance and safety mature over time
      \item Backward compatibility is mandatory
    \end{itemize}
  }
  {
    \focusline{Evolution must preserve operational stability.}
  }

\twocolslide
  {}
  {SEMEIA Evolution}
  {Roadmap (Evolution)}
  {
    \begin{itemize}
      \item \textbf{Now}: roadmap governance cadence + KPI baseline
      \item \textbf{Next}: expand pilot workflows + runbooks
      \item \textbf{Later}: scale integrations + institutionalize evaluation
    \end{itemize}
  }
  {
    \begin{itemize}
      \item Owners: Tech Lead + Product + Ops
      \item Risks: demand growth > capacity; ops burden
    \end{itemize}
  }

% -----------------------------------------------------
% Final slide – Thank you
% -----------------------------------------------------
\begin{frame}[plain]
  \frameref{}%
  \begin{columns}[T,onlytextwidth]
    \begin{column}{0.62\textwidth}
      \vspace{0.85cm}
      {\LARGE\bfseries \DeckTitle\par}

      \vspace{0.95cm}
      {\Large\bfseries Thank you!}

      \vspace{0.95cm}
      {\small
        \texttt{\PresenterName}\\[0.15em]
        \texttt{\PresenterEmail}\\[0.15em]
        \texttt{\OrgLine}\\[0.15em]
      }
    \end{column}

    \begin{column}{0.38\textwidth}
      \vspace{0.20cm}
      \hfill
      \deckgraphic[height=0.86\textheight]{\logoBig}
    \end{column}
  \end{columns}
\end{frame}

\end{document}
